% Shared preamble for the CubeSat coupling writeups.
% Build any writeup with: pdflatex NAME; bibtex NAME; pdflatex NAME; pdflatex NAME
\documentclass[11pt]{article}
\usepackage[margin=1in]{geometry}
\usepackage[T1]{fontenc}
\usepackage{lmodern}
\usepackage{amsmath,amssymb}
\usepackage{siunitx}
\usepackage{booktabs}
\usepackage{enumitem}
\usepackage{xcolor}
\usepackage[numbers,sort&compress]{natbib}
\usepackage[colorlinks=true,linkcolor=blue!50!black,citecolor=blue!50!black,urlcolor=blue!50!black]{hyperref}

\sisetup{detect-all, per-mode=symbol}
\setlist{nosep, leftmargin=*}

% Code paths, relative to the repository roots
\newcommand{\code}[1]{\texttt{\detokenize{#1}}}
% MC/DC repository (feature/g4-hpc) and Geant4 bridge repository (main)
\newcommand{\mcdc}{MC/DC}
\newcommand{\gfour}{Geant4}

% A labelled assumption, numbered within each writeup
\newcounter{assumption}
\newcommand{\assumption}[1]{\refstepcounter{assumption}\par\smallskip\noindent\textbf{Assumption A\theassumption.}~#1\par\smallskip}

\bibliographystyle{unsrtnat}


\title{CubeSat coupling writeup 6:\\Energy-dependent reaction yields (ACE $|TY|>100$) and the data patch}
\author{CARRE CubeSat MC/DC--Geant4 coupling}
\date{\today}

\begin{document}
\maketitle

\begin{abstract}
Runs before 21738977 overproduced neutrons by a large factor because the neutron data library stored an encoded yield flag as a literal multiplicity. This note describes the error, the code fix that was ported to \code{feature/g4-hpc}, the in-place patch of the combined neutron--proton library used by the CubeSat runs, and how the fix was verified.
\end{abstract}

\section{The ACE yield convention}
In an ACE continuous-energy table, each reaction's TY value gives the number of secondary neutrons and the frame of their distributions. $|TY|$ of 1--4 is a fixed integer multiplicity; $|TY|>100$ means the multiplicity is energy dependent and tabulated, with $|TY|-100$ an offset into the DLW block where a yield table $\nu(E)$ is stored; the sign gives the reference frame \citep{conlin2019ace}.

MT-5 is the ENDF reaction $(z,\text{anything})$, the sum of all reactions not given explicitly \citep{trkov2018endf6}. In the JENDL-5 \citep{iwamoto2023jendl5} ACE files of the CubeSat nuclides (and the FENDL-3.0 file for Li-7), MT-5 has an energy-dependent neutron yield, written as $TY=101$ (Table~\ref{tab:patched}).

\section{The error}
The \mcdc{} neutron data generator (\code{tools/data_library_generator/neutron}) wrote the raw TY value as an integer multiplicity. The transport then emitted 101 neutrons in every MT-5 event. In the CubeSat, neutrons therefore multiplied in every structure: the neutron current into the electronics boxes was 5 to 10 times the unshielded value (run 21616157; Table~\ref{tab:verify}), and neutron-source energy deposition was inflated accordingly.

\section{Code fix}
Commit \code{0751e69a} (ported to \code{feature/g4-hpc}):
\begin{itemize}
  \item the generator writes, for $|TY|>100$, a group \code{multiplicity_table} with the yield $\nu(E)$ (lin-lin, $E$ in MeV) and sets \code{multiplicity = -1};
  \item the inelastic reaction object carries the table, and warns if old data still carries an encoded multiplicity;
  \item the sampler emits $n=\lfloor \nu(E) + \xi\rfloor$ secondaries, $\xi\sim U(0,1)$. With $\nu = m + f$ ($m$ integer, $0\le f<1$), $n=m+1$ with probability $f$ and $m$ otherwise, so $\mathbb{E}[n]=\nu(E)$.
\end{itemize}
Commit \code{be31c7b5} makes the unit test \code{test/unit/transport/test_tabulated_multiplicity.py} run with neutron transport enabled, so it exercises the sampler.

\section{Patching the library in place}
The combined library \code{hdf5libProton/with_neutron_200MeV} (JENDL-5 neutrons, FENDL-3.0 \citep{forrest2012fendl3} for Li-7, TENDL-2021 \citep{koning2019tendl} protons) was patched rather than regenerated, with the script \code{patch_tabulated_yields.py}. For each nuclide it loads the ACE file the library was generated from (JENDL-5 ACE, \code{ace_j50_neu_0300.0k} \citep{jendl5ace}; the FENDL-3.0 file for Li-7) and, before touching a reaction, checks that:
\begin{enumerate}
  \item the HDF5 energy grid equals the ACE principal energy grid;
  \item the reaction's stored multiplicity has $|TY|>100$ and has no yield table yet;
  \item the ACE file has the same MT with $|TY|>100$;
  \item the stored cross section and its grid offset equal the ACE values exactly.
\end{enumerate}
Only then is the multiplicity replaced by $-1$ and the ACE yield table written, with the same routine the fixed generator uses. Any failed check aborts. The unpatched library is kept in \code{with_neutron_200MeV_preMT5fix}, and the HPC copy was verified identical to the local patched copy by MD5 checksum.

\begin{table}[h]
\centering
\caption{Reactions patched. All are MT-5 with $TY=101$; tables span \SI{1e-11}{MeV} to \SI{200}{MeV}.}
\label{tab:patched}
\begin{tabular}{lrl}
\toprule
Nuclide & Points & Yield range $\nu(E)$ \\
\midrule
C-12  & 25 & 0--1.85 \\
Li-7  & 42 & 0.127--1.03 \\
O-16  & 33 & 0--1.73 \\
Mg-24 & 29 & 0--2.15 \\
Al-27 & 33 & 0--2.48 \\
Si-28 & 63 & 0--1.94 \\
Co-59 & 43 & 0--3.1 \\
Cu-63 & 32 & 0--2.91 \\
Cu-65 & 32 & 0--3.42 \\
Zn-64 & 43 & 0--2.59 \\
\bottomrule
\end{tabular}
\end{table}

\section{Verification}
The CubeSat walls and boards are thin (of order \si{g/cm^2}) and neutrons are rarely absorbed or multiplied in them, so the neutron current into the boxes is expected to stay close to the unshielded value: scattering mainly redistributes neutrons in direction and energy. A transmission of 5--10 requires a large net source of neutrons inside the spacecraft. We use $T\approx1$ as a consistency check, not as an exact expectation; capture lowers it and $(n,xn)$ reactions and proton-born neutrons raise it. Table~\ref{tab:verify} compares runs before and after the patch.

\begin{table}[h]
\centering
\caption{Neutron transmission into each box, $T=$ crossings / unshielded crossings (writeup 5), with \mcdc{} standard errors.}
\label{tab:verify}
\begin{tabular}{lccc}
\toprule
Region & 21616157 (before) & 21738977 (after, $N_p=10^6$) & 21756183 (after, $N_p=10^7$) \\
\midrule
OBC   & $7.22\pm0.02$ & $1.004\pm0.007$ & $1.021\pm0.007$ \\
EPS   & $7.54\pm0.02$ & $1.001\pm0.006$ & $1.022\pm0.006$ \\
ADCS  & $5.12\pm0.01$ & $0.988\pm0.004$ & $1.017\pm0.004$ \\
COMMS & $9.64\pm0.06$ & $0.998\pm0.020$ & $0.988\pm0.020$ \\
\bottomrule
\end{tabular}
\end{table}

After the patch the transmission is $1.00\pm0.02$, and the neutron-source event rates in the \gfour{} stage fell by a factor of 5--6 at all charges. The small rise from $N_p=10^6$ to $10^7$ is the proton-born neutron contribution (writeup 2, species attribution).

\assumption{Results from runs before 21738977 (including 21616157) have the wrong neutron transport and are not used for neutron-source quantities.}

\bibliography{refs}
\end{document}
